\documentclass[../../main.tex]{subfiles}
\begin{document}


\begin{problem}
    n=7,M=15,

    \((P_1,P_2,P_3,P_4,P_5,P_6,P_7)=(10,5,15,7,6,18,3)\)

    \((w_1,w_2,w_3,w_4,w_5,w_6,w_7)=(2,3,5,7,1,4,1)\)

    将以上数据情况的背包问题记为I。
    
    设\(FG(I)\)是物品按照\(P_i\)的非增次序输入时由GREEDY-KNAPSACK所生成的解,\(FO(I)\)是一个最优解。
\begin{itemize}
    \item 问FO(I)/FG(I)是多少？
    \item 当物品按照\(w_i\)的非降次序输入时,重复上述讨论。
\end{itemize}
\end{problem}
\begin{solution}
\begin{itemize}
    \item 问FO(I)/FG(I)是多少？

首先按物品的价值排序。物品 \( P \) 的非增次序排序为:
\begin{itemize}
    \item 物品 6 (价值 \( P_6 = 18 \), 重量 \( w_6 = 4 \), 单位价值 \(4.5 \))
    \item 物品 3 (价值 \( P_3 = 15 \), 重量 \( w_3 = 5 \), 单位价值 \(3   \))
    \item 物品 1 (价值 \( P_1 = 10 \), 重量 \( w_1 = 2 \), 单位价值 \(5   \))
    \item 物品 4 (价值 \( P_4 =  7 \), 重量 \( w_4 = 7 \), 单位价值 \(1   \))
    \item 物品 5 (价值 \( P_5 =  6 \), 重量 \( w_5 = 1 \), 单位价值 \(6   \))
    \item 物品 2 (价值 \( P_2 =  5 \), 重量 \( w_2 = 3 \), 单位价值 \(1.66\))
    \item 物品 7 (价值 \( P_7 =  3 \), 重量 \( w_7 = 1 \), 单位价值 \(3   \))
\end{itemize}
计算\(FG(I)\):
\begin{itemize}
\item[(1)] 物品 6:价值 18,重量 4,背包剩余容量 \( 15 - 4 = 11 \),选择物品 6。
\item[(2)] 物品 3:价值 15,重量 5,背包剩余容量 \( 11 - 5 = 6 \),选择物品 3。
\item[(3)] 物品 1:价值 10,重量 2,背包剩余容量 \( 6 - 2 = 4 \),选择物品 1。
\item[(4)] 物品 4:价值 7,重量 7,背包剩余容量 \( 4 - 7 = -3 \),只能放入\(\frac{4}{7}\)。
\end{itemize}
对应的解向量为 \( X = (1, 0, 1, \frac{4}{7}, 0, 1, 0) \),目标值为:
\[
FG(I) = \sum P_i x_i = 18 + 15 + 10 + 4= 47
\]

计算\( FO(I) \):
按物品的单位价值排序:
\begin{itemize}
    \item 物品 5 (价值 \( P_5 =  6 \), 重量 \( w_5 = 1 \), 单位价值 \(6   \))
    \item 物品 1 (价值 \( P_1 = 10 \), 重量 \( w_1 = 2 \), 单位价值 \(5   \))
    \item 物品 6 (价值 \( P_6 = 18 \), 重量 \( w_6 = 4 \), 单位价值 \(4.5 \))
    \item 物品 3 (价值 \( P_3 = 15 \), 重量 \( w_3 = 5 \), 单位价值 \(3   \))  
    \item 物品 7 (价值 \( P_7 =  3 \), 重量 \( w_7 = 1 \), 单位价值 \(3   \))

    \item 物品 2 (价值 \( P_2 =  5 \), 重量 \( w_2 = 3 \), 单位价值 \(1.66\))
    \item 物品 4 (价值 \( P_4 =  7 \), 重量 \( w_4 = 7 \), 单位价值 \(1   \))
\end{itemize}
那么解向量为 \( X = (1, \frac{2}{3}, 1, 0, 1, 1, 1) \),目标值为:
\[
FO(I) = \sum P_i x_i =  6+10+18+15+3+5 \times \frac{2}{3}  = 55.33
\]
\(FO(I)/FG(I) = 55.33/47 = 1.18\)

\item 当物品按照\(w_i\)的非降次序输入时,重复上述讨论。

按物品的重量排序:
\begin{itemize}
    \item 物品 5 (价值 \( P_5 =  6 \), 重量 \( w_5 = 1 \), 单位价值 \(6   \))
    \item 物品 7 (价值 \( P_7 =  3 \), 重量 \( w_7 = 1 \), 单位价值 \(1   \))
    \item 物品 1 (价值 \( P_1 = 10 \), 重量 \( w_1 = 2 \), 单位价值 \(5   \))
    \item 物品 2 (价值 \( P_2 =  5 \), 重量 \( w_2 = 3 \), 单位价值 \(1.66\))
    \item 物品 6 (价值 \( P_6 = 18 \), 重量 \( w_6 = 4 \), 单位价值 \(4.5 \))
    \item 物品 3 (价值 \( P_3 = 15 \), 重量 \( w_3 = 5 \), 单位价值 \(3   \))
    \item 物品 4 (价值 \( P_4 =  7 \), 重量 \( w_4 = 7 \), 单位价值 \(1   \))
\end{itemize}
那么解向量为 \( X = (1, 1, \frac{4}{5}, 0, 1, 1, 1) \),目标值为:
\[
FG(I) = \sum P_i x_i =  6+3+10+5+18+15 \times \frac{4}{5}  = 54
\]
\(FO(I)/FG(I) = 55.33/54 = 1.02\)
\end{itemize}
\end{solution}
\end{document}